\documentclass[letterpaper,11pt]{moderncv}
\moderncvstyle{banking}
\moderncvcolor{black}

\usepackage[letterpaper, margin=0.8in]{geometry}
\usepackage{xcolor}
\usepackage{fontspec}
\setmainfont{Times New Roman}
\usepackage[unicode]{hyperref}

\name{Saeyoung}{Kim}
\address{Cambridge, MA, USA}{}{}
\phone[mobile]{+1 (351) 322 5510}
\email{saeyoung\_kim@uml.edu}
\social[linkedin]{saeyoung-macx-kim}
\social[github]{PrimaryAnomaly}

\recipient{Hiring Team}{Boston Dynamics\\Waltham, MA}
\date{\today}
\opening{Dear Hiring Team,}
\closing{Sincerely,}

\begin{document}

\makelettertitle

I am applying for the Robotics Technician V position (R2277) on the Spot Engineering team. I have spent the last 9+ years building, testing, and debugging dense electromechanical systems, and I want to do that work on robots.

At Hanwha Systems, I built and integrated satellite sub-assemblies: computers, power distribution, embedded processors. I wrote my own test plans, ran design validation on new subsystems, and tracked down root causes when things failed. I did most of this independently. I also designed test fixtures, set up lab spaces, and kept long-running test campaigns on track.

During my PhD, I built prototype electromechanical systems from scratch. PCB design, layout, assembly, sensor integration, firmware in C/C++, system-level validation. I did PCB rework and cable assembly daily, and I lived on oscilloscopes, power supplies, and thermal probes. Now in my postdoc, I write test automation in Python on Linux and manage long-running test-stand experiments.

I want to work at Boston Dynamics because I do my best work when I am building and troubleshooting hardware in a fast-moving R\&D environment. I am used to juggling multiple tasks, shifting priorities quickly, and delivering clear test results. I would love to bring that to the Spot team.

Thank you for your consideration.

\makeletterclosing

\end{document}
